\documentclass[../../main.tex]{subfiles}
\begin{document}

{\bf[解]}

簡単のため，$f_n(x),g_n(x)$の定義域を$n\ge 0$に拡張し
\begin{align}
  \begin{cases}
    f_{0}(x) = f(x) \\
    g_{n}(c) = \int_{0}^{c} f_n(x) dx
  \end{cases} \label{eq:1}
\end{align}
である．
与えられた漸化式
\begin{align}
  \begin{cases}
    f_{n+1}(x) = f(x) + \int_{0}^{c} f_{n}(x) dx & \left(n=0,1,2,\cdots \right)\\
    g_{n}(c) = \int_{0}^{c} f_{n}(x) dx & \left(n=0,1,2,\cdots \right)
  \end{cases} \label{eq:2}
\end{align}
について考える．

\cref{eq:2}の第二式を第一式に代入して
\begin{align}
  f_{n+1}(x) = f(x) + g_{n}(c) \label{eq:3}
\end{align}
だから，
\begin{align}
  g_{n+1}(x)
  & = \int_{0}^{c} \left[f(x) + g_{n}(c)\right] dx \\
  & = g(x) + cg_{n}(c)
\end{align}
である．$0<c<1$に注意して変形すると
\begin{align}
  g_{n+1}(c)-\dfrac{g(c)}{1-c}&=c\left\{ g_{n}(c)-\dfrac{g(c)}{1-c}\right\} \label{eq:4}
\end{align}
である．
これは数列$\{ g_n(c)-\dfrac{g(c)}{1-c}\}$が公比$c$の等比数列であることをあらわす．
また，初期条件は$g_{0}(c)=g(c)$だから
\begin{align}
  g_n(c)-\dfrac{g(c)}{1-c} &= c^{n}\left[g(c)-\dfrac{g(c)}{1-c}\right] \\
  \therefore
  g_n(c)
  &= c^{n}\left[g(c)-\dfrac{g(c)}{1-c}\right] + \dfrac{g(c)}{1-c} \\
  &= \dfrac{-cg(c)}{1-c}c^{n} + \dfrac{g(c)}{1-c} \\
  &= \dfrac{1-c^{n+1}}{1-c}g(c) \label{eq:5}
\end{align}
である．

従って$c$を$x$で置き換えて$0<x<1$より，
\begin{align}
  \lim_{n\to\infty}g_{n}(x) = \dfrac{g(x)}{1-x} \label{eq:6}
\end{align}
だから題意の方程式
\begin{align}
  xf(x) = g(x) + x\lim_{n\to\infty}g_{n}(x) \label{eq:7}
\end{align}
に代入して
\begin{align}
  xf(x) &= g(x) + \dfrac{xg(x)}{1-x} \\
  (1-x)x f(x) &= g(x) \label{eq:8}
\end{align}
である．以下これを満たす$f(x)$を求めよう．
まず$y=g(x)$とおくと，\cref{eq:1}より$f(x)=\dfrac{dy}{dx}$だから\cref{eq:8}より
\begin{align}
  y &= (1-x)x \dfrac{dy}{dx} \\
  \dfrac{dy}{y} &= \dfrac{dx}{x(1-x)} = \left(\dfrac{1}{x}+\dfrac{1}{1-x}\right)dx
\end{align}
である．両辺積分して$C$を積分定数として
\begin{align}
  \log y
  &= \log x - \log (1-x) + C \\
  &= \log \dfrac{x}{1-x} + C \\
  \therefore
  y&=e^{C}\dfrac{1}{1-x} \quad (0<x<1) \label{eq:9}
\end{align}
である．$y=f(x)$の連続性から $x=0$ でも\cref{eq:9}で良い．
題意より$f(0)=e^{C}=1$だから
\begin{align}
  f(x) = \dfrac{1}{(1-x)^2} \quad (0<x<1) \label{eq:10}
\end{align}
である．

\newpage
\end{document}
